\scanpage{140}
% lines 1–6
2) Soit $S'\unsure{-|\to}S$ fid. plat et qu. cpt. Alors

(i) $\mathrm{Pic}'(X/S)-\mathrm{Pic}'(X'/S')$ est injectif. (ii) Si $X/\Sigma$ a \deleted{\unsure{loc.}} une section, \deleted{\unsure{sur}} localement sur $S$, alors
\[
\mathrm{Pic}'(X/S)\to\mathrm{Pic}'(X'/S')\rightrightarrows\mathrm{Pic}'(X''/S'')
\]
est exact.
% slanted
\margin{Généraliser à des ``quasi-sections'' du type \unsure{$N^{*}$|$N^{\circ}$} dès l'abord \textemdash}

% lines 7–10
On peut dans (i) partir d'un $\underline L$ inv. sur $X$. \inserted{En vertu de 1),} L'hyp. signifie que $\underline L'$ provient de $\Sigma'$. Or $\Sigma'$ est fid plat qu. cpt sur $\Sigma$. On sait qu'alors \deleted{\unsure{le}} $L$ provient de $\Sigma$.

% lines 11–15
Pour (ii), \inserted{\unsure{grâce à|après|qui est} (i),} la question est locale sur $S$, on peut donc supposer que $X$ a une section sur $\Sigma$. Alors on utilise la théorie de descente comme dans le raisonnement connu.

% lines 16–20
\uline{Corollaire} \deleted{\unsure{1}} Sous les conditions précédentes, $\mathrm{Pic}'(X/S)\to\mathrm{Pic}^{*}(X/S)$ est injectif, et \unsure{même|est} bijectif si $X/\Sigma$ \deleted{\unsure{loc une}} a une section loc. sur $S$ ; Si $X/\Sigma$ a \unsure{\textnormal{même}} une section, on a un isom
\[
\mathrm{Pic}(X/S)\simeq\mathrm{Pic}(X)/\mathrm{Pic}(\Sigma)
\]
% archive sheet number: 139

